\chapter{Clientes y usuarios}
\label{cap:clientes}

En BeeTracer existen dos grandes grupos de personas según el rol que cumplen dentro del sistema:
por una parte, los \textbf{trabajadores de la empresa de logística} que operan la plataforma web o
la aplicación móvil y, por otra, los \textbf{clientes}, que esperan información sobre la mercancía
que tienen en el depósito. Cada grupo cuenta con distintos perfiles y niveles de autorización, lo
que permite organizar de mejor forma las funcionalidades y responsabilidades. Todos los perfiles,
con excepción del Servicio Nacional de Aduanas, deben iniciar sesión con sus credenciales para
acceder al sistema con los permisos de su rol.

\section{Usuarios trabajadores de la empresa de logística}
\label{sec:usuarios_trabajadores}

\textbf{Administrador del depósito:} es el usuario con el nivel más alto de privilegios. Configura
los usuarios, roles, clientes, trabajadores, cuadrillas y tipos de bulto de la plataforma, y define
quién puede acceder a qué funcionalidades.

\textbf{Jefe de faena:} actúa como planificador de las operaciones. Desde la plataforma web importa
los manifiestos, programa los desconsolidados, decide qué cuadrilla realizará cada faena y quién
será el jefe tarjador, y supervisa el avance de las operaciones: las consulta, modifica su estado,
visualiza o descarga su reporte PDF y, una vez aprobado, lo envía al cliente.

\textbf{Cuadrilla:} conjunto de cuatro a cinco trabajadores encargados de manipular físicamente los
bultos del contenedor, separándolos por cliente y apartando los dañados. No opera el sistema
directamente: su trabajo se registra a través del jefe tarjador, por lo que no se modela como
actor.

\textbf{Jefe tarjador:} integrante de la cuadrilla que opera la aplicación móvil en terreno. Inicia
y finaliza la faena y registra el sello, las cantidades, las incidencias y la evidencia
fotográfica. Es el usuario más crítico del sistema, ya que con su registro se genera el reporte
para cada cliente; trabaja con guantes y a la intemperie, por lo que requiere una interfaz simple
y guiada.

\textbf{Supervisor del reporte:} usuario interno que valida los reportes generados (datos y
fotografías) antes de que lleguen al cliente. Verifica que la evidencia sea consistente y legible,
marca la evidencia defectuosa, rechaza el reporte registrando una observación y solicitando la
corrección, o lo aprueba para que pueda ser enviado.

\section{Clientes}
\label{sec:clientes}

\textbf{Empresa dueña del depósito:} es el cliente del sistema, es decir, la organización
(terminal extraportuario o bodega) que contrata BeeTracer bajo suscripción para mejorar la
eficiencia de su operación y protegerse ante reclamos. Sus trabajadores son los usuarios descritos
en la sección anterior.

\textbf{Cliente del depósito:} son las empresas importadoras (zapaterías, jugueterías,
ferreterías) cuya carga viene consignada en el contenedor. No configuran ni operan el sistema: su
interacción es de consulta, visualizan el estado de sus operaciones y visualizan o descargan los
reportes aprobados. Como no acceden a la evidencia en vivo, el terminal puede revisar la
información antes de notificarla.

\textbf{Servicio Nacional de Aduanas:} es un actor externo (sistema). Pone a disposición el
manifiesto electrónico en formato XML que BeeTracer descarga para conocer el contenido de los
contenedores antes de su arribo, y recibe la notificación del resultado de la importación. No
inicia sesión en el sistema.
